\documentclass[a4paper]{article}
\usepackage{vntex}
\usepackage[T5]{fontenc}
\usepackage[utf8]{inputenc}
\usepackage{lmodern}
\usepackage{pdflscape}
\usepackage{rotating}
\usepackage{a4wide,amssymb,epsfig,latexsym,multicol,array,hhline,fancyhdr}
\usepackage{caption}
\usepackage{subcaption}
\usepackage{booktabs}
\usepackage{amsmath}
\usepackage{float}
\usepackage{placeins}
\usepackage{lastpage} 
% \usepackage[lined,boxed,commentsnumbered]{algorithm2e}
\usepackage{enumerate}
\usepackage{color}
\usepackage{graphicx}					% Standard graphics package
\usepackage{array}
\usepackage{tabularx, caption}
\usepackage{multirow}
\usepackage[framemethod=tikz]{mdframed}% For highlighting paragraph backgrounds
\usepackage{multicol}
\usepackage{rotating}
\usepackage{graphics}
\usepackage{wrapfig}
\usepackage{geometry}
\usepackage{setspace}
\usepackage{epsfig}
\usepackage{tikz}
\usepackage{listings}
\usetikzlibrary{arrows,snakes,backgrounds}
\usepackage{hyperref}
\usepackage{indentfirst}
\usepackage{cases} 
\usepackage{listings}
\usepackage{algorithm}
\usepackage{algpseudocode}
% \usepackage[english]{babel}
\renewcommand{\baselinestretch}{1.5}
\usepackage{xcolor}
\renewcommand{\figurename}{Hình}
\usepackage{xcolor}
\definecolor{light-gray}{gray}{0.95}
\newcommand{\code}[1]{\colorbox{light-gray}{\texttt{#1}}}

\lstset{
    language=C++,
    basicstyle=\ttfamily\footnotesize,
    numbers=left,
    numberstyle=\tiny,
    stepnumber=1,
    numbersep=5pt,
    backgroundcolor=\color{gray!10},
    keywordstyle=\color{blue}\bfseries,
    commentstyle=\color{green!50!black},
    stringstyle=\color{red},
    frame=single,
    breaklines=true,
}

%%%%%%%%
\usepackage{longtable}
\usepackage[shortlabels]{enumitem}
\hypersetup{urlcolor=blue,linkcolor=black,citecolor=black,colorlinks=true} 

\setlength{\headheight}{40pt}
\pagestyle{fancy}
\fancyhead{} % clear all header fields
\fancyhead[L]{
 \begin{tabular}{rl}
    \begin{picture}(25,15)(0,0)
    \put(0,-8){
\includegraphics[width=8mm, height=8mm]{images/khtn.png}}
    \end{picture}&
	\begin{tabular}{l}
		\textbf{\bf \ttfamily Trường đại học Khoa Học Tự Nhiên TP.HCM}\\
		\textbf{\bf \ttfamily Khoa Công Nghệ Thông Tin}
	\end{tabular}  	
 \end{tabular}
}
\fancyhead[R]{
	\begin{tabular}{l}
		\tiny \bf \\
		\tiny \bf 
	\end{tabular}  }
\fancyfoot{} % clear all footer fields
\fancyfoot[L]{\scriptsize \ttfamily Đồ án 2 | Data Fitting và Phương Pháp OLS }
\fancyfoot[R]{\scriptsize \ttfamily Trang {\thepage}}
\renewcommand{\headrulewidth}{0.3pt}
\renewcommand{\footrulewidth}{0.3pt}

\fancypagestyle{plain}{
\fancyhead{}
\fancyhead[L]{
 \begin{tabular}{rl}
    \begin{picture}(25,15)(0,0)
    \put(0,-8){
\includegraphics[width=8mm, height=8mm]{images/khtn.png}}
    \end{picture}&
    \begin{tabular}{l}
        \textbf{\bf \ttfamily Trường đại học Khoa Học Tự Nhiên TP.HCM}\\
        \textbf{\bf \ttfamily Khoa Công Nghệ Thông Tin}
    \end{tabular}
 \end{tabular}
}
\fancyhead[R]{
    \begin{tabular}{l}
        \tiny \bf \\
        \tiny \bf
    \end{tabular}  }
\fancyfoot{}
\fancyfoot[L]{\scriptsize \ttfamily Đồ án 2 | Data Fitting và Phương Pháp OLS}
\fancyfoot[R]{\scriptsize \ttfamily Trang {\thepage}}
\renewcommand{\headrulewidth}{0.3pt}
\renewcommand{\footrulewidth}{0.3pt}
}


%%%
\setcounter{secnumdepth}{4}
\setcounter{tocdepth}{3}
\makeatletter
\newcounter {subsubsubsection}[subsubsection]
\renewcommand\thesubsubsubsection{\thesubsubsection .\@alph\c@subsubsubsection}
\newcommand\subsubsubsection{\@startsection{subsubsubsection}{4}{\z@}%
                                     {-3.25ex\@plus -1ex \@minus -.2ex}%
                                     {1.5ex \@plus .2ex}%
                                     {\normalfont\normalsize\bfseries}}
\newcommand*\l@subsubsubsection{\@dottedtocline{3}{10.0em}{4.1em}}
\newcommand*{\subsubsubsectionmark}[1]{}
\makeatother

\definecolor{dkgreen}{rgb}{0,0.6,0}
\definecolor{gray}{rgb}{0.5,0.5,0.5}
\definecolor{mauve}{rgb}{0.58,0,0.82}

\lstset{frame=tb,
	language=Matlab,
	aboveskip=3mm,
	belowskip=3mm,
	showstringspaces=false,
	columns=flexible,
	basicstyle={\small\ttfamily},
	numbers=none,
	numberstyle=\tiny\color{gray},
	keywordstyle=\color{blue},
	commentstyle=\color{dkgreen},
	stringstyle=\color{mauve},
	breaklines=true,
	breakatwhitespace=true,
	tabsize=3,
	numbers=left,
	stepnumber=1,
	numbersep=1pt,    
	firstnumber=1,
	numberfirstline=true
}

\begin{document}

% --- PHẦN NỘI DUNG TỪ title.tex ---
\begin{titlepage}
\begin{center}
ĐẠI HỌC QUỐC GIA THÀNH PHỐ HỒ CHÍ MINH\\
TRƯỜNG ĐẠI HỌC KHOA HỌC TỰ NHIÊN\\
KHOA CÔNG NGHỆ THÔNG TIN
\end{center}

\vspace{0.7cm}

\begin{figure}[h!]
\begin{center}

\includegraphics[width=5cm]{images/khtn.png}
\end{center}
\end{figure}

\vspace{0.5cm}


\begin{center}
\textbf{\large TOÁN ỨNG DỤNG VÀ THỐNG KÊ}\\[0.2cm]
\rule{0.8\textwidth}{0.4pt}\\[0.3cm]
\textbf{\Large Đồ án 02}\\[0.2cm]
\textbf{\Large Data Fitting và Phương Pháp OLS}\\[0.3cm]
\rule{0.8\textwidth}{0.4pt}\\[0.2cm]
\textbf{\large THỰC HIỆN BỞI NHÓM 1}
\end{center}

\vspace{0.1cm}

\begin{center}
\begin{tabular}{rll}
\hspace{1.8cm} & \textbf{Sinh viên thực hiện:} & Trần Nhật Trường - 24120486  \\
                                      &                    & Nguyễn Thanh Nhật - 24120111 \\
                                      &                    & Nguyễn Ngọc Phúc - 24120215 \\
                                      &                    & Trần Lê Đức Việt - 24120245 \\
                                      &                    & Trần Nguyên Tân - 24120438 \\
\hspace{1.8cm} & \textbf{Giảng viên hướng dẫn:} & ThS. Lê Nhựt Nam \\
                                            &                &    ThS. Võ Nam Thục Đoan \\    


\end{tabular}
\end{center}

\vfill

\begin{center} {\footnotesize Thành phố Hồ Chí Minh, 05/2026}
\end{center}
\end{titlepage}

\clearpage
\pagestyle{fancy}
\pagenumbering{roman}
\renewcommand{\contentsname}{Mục lục}
\renewcommand{\tablename}{Bảng}
\renewcommand{\figurename}{Hình}
\renewcommand{\refname}{Tài liệu tham khảo}

\tableofcontents
\thispagestyle{fancy}

\newpage

\renewcommand{\arraystretch}{1.3} % Adjust row height
\pagenumbering{arabic}  

% --- PHẦN NỘI DUNG TỪ sample.tex ---

\section*{Lời nói đầu}
\addcontentsline{toc}{section}{Lời nói đầu}

Lời đầu tiên, tất cả thành viên nhóm gửi lời cảm ơn chân thành đến thầy Lê Nhựt Nam và cô Võ Nam Thục Đoan đã tận tình hướng dẫn, hỗ trợ chúng em trong suốt quá trình thực hiện đồ án này. Trong suốt quá trình làm đồ án này, với sự hướng dẫn tận tình của thầy và cô đã giúp chúng em có cái nhìn tổng quan, sâu sắc và rõ hơn về đồ án ma trận và cơ sở của tính toán khoa học. Chúng em đã học được rất nhiều kiến thức mới, cũng như kỹ năng làm việc nhóm, giải quyết vấn đề và trình bày kết quả một cách hiệu quả.

Chúng em cũng xin gửi sự trân trọng và biết ơn đến các thành viên trong nhóm thực hiện đồ án này đã làm việc chăm chỉ và năng nổ trong suốt quá trình thực hiện. Sự đóng góp to lớn của mỗi thành viên đã góp phần quan trọng vào thành công của đồ án này. Chúng em đã học được rất nhiều từ nhau, cũng như phát triển kỹ năng làm việc nhóm và giao tiếp hiệu quả.    

Trong lúc thực hiện đồ án này không tránh khỏi những khó khăn và thách thức nên sẽ có những sai sót không mong muốn. Chúng em rất mong nhận được sự góp ý, phản hồi từ thầy và cô để có thể hoàn thiện hơn trong quá trình học tập và nghiên cứu sau này.

\section*{Lời cam đoan}
\addcontentsline{toc}{section}{Lời cam đoan}

Trong đồ án lần này, chúng em cam đoan rằng tất cả các nội dung mà chúng em thực hiện đều là kết quả của sự nghiên cứu và làm việc của nhóm chung em. 

Chúng em cam đoan rằng không sao chép code hoặc báo cáo từ nhóm khác hay từ Internet mà không
trích dẫn nguồn. 

Chúng em cam kết khi sử dụng các công cụ AI (ChatGPT, GitHub Copilot, v.v.) để gợi ý thì phải hiểu, trình bày và giải thích lại được nội dung đã được gợi ý.

Chúng em cam đoan rằng tất cả các thành viên trong nhóm đều đã đóng góp công sức và kiến thức của mình vào đồ án này một cách công bằng và trung thực. Chúng em cam đoan rằng nếu có bất kỳ vi phạm nào liên quan đến đạo đức học thuật hoặc gian lận, chúng em sẽ chịu trách nhiệm hoàn toàn và chấp nhận mọi hình phạt theo quy định của nhà trường.

\section{Bảng phân công công việc}
% \addcontentsline{toc}{section}{Bảng phân công công việc}

\begingroup
\small
\setlength{\tabcolsep}{4pt}
\renewcommand{\arraystretch}{1.15}
\begin{longtable}{|>{\raggedright\arraybackslash}p{0.20\textwidth}|>{\raggedright\arraybackslash}p{0.62\textwidth}|>{\centering\arraybackslash}p{0.10\textwidth}|}

\hline
\textbf{Tên} & \textbf{Nội dung} & \textbf{Phần trăm đánh giá} \\ \hline
\endfirsthead
\hline
\textbf{Tên} & \textbf{Nội dung} & \textbf{Phần trăm đánh giá} \\ \hline
\endhead
Nguyễn Thanh Nhật   & \begin{itemize}[leftmargin=*, noitemsep, topsep=0pt, partopsep=0pt]
                        \item Trình bày và chứng minh lý thuyết OLS
                        \item Cài đặt \texttt{ols\_fit} và \texttt{hat\_matrix}
                        \item Viết unit test cho OLS và hat matrix
                        \item Minh họa định lý Gauss--Markov bằng Monte Carlo
                        \item Bổ sung Markdown phần công thức cốt lõi trong notebook
                    \end{itemize} & 100\% \\ \hline
Trần Lê Đức Việt   & \begin{itemize}[leftmargin=*, noitemsep, topsep=0pt, partopsep=0pt]
                        \item Cài đặt \texttt{model\_metrics}, \texttt{coef\_inference}, \texttt{vif}
                        \item Viết unit test cho các hàm kiểm định thống kê
                        \item Bổ sung Markdown giải thích F-test, t-test, RSS, $R^2$
                    \end{itemize} & 100\% \\ \hline
Trần Nhật Trường & \begin{itemize}[leftmargin=*, noitemsep, topsep=0pt, partopsep=0pt]
                        \item Thực hiện EDA dữ liệu Air Quality
                        \item Phân tích cơ chế missing values và outliers
                        \item Thiết kế \texttt{DataPipeline} (fit/transform, imputation, encoding, Z-score)
                        \item Bổ sung Markdown mô tả EDA và pipeline trong notebook
                    \end{itemize} & 100\% \\ \hline
Trần Nguyên Tân   & \begin{itemize}[leftmargin=*, noitemsep, topsep=0pt, partopsep=0pt]
                        \item Cài đặt Ridge Regression và k-fold CV
                        \item Vẽ Ridge Trace và chọn lambda tối ưu
                        \item Tích hợp so sánh mô hình trong \texttt{model\_comparison.py}
                        \item Triển khai Bayesian Linear Regression (bonus)
                        \item Viết nhận xét kết quả so sánh mô hình
                    \end{itemize} & 100\% \\ \hline
Nguyễn Ngọc Phúc     & \begin{itemize}[leftmargin=*, noitemsep, topsep=0pt, partopsep=0pt]
                        \item Cài đặt \texttt{residual\_plots} và unit test
                        \item Vẽ và nhận xét 4 biểu đồ phần dư cho Phần 1
                        \item Vẽ Feature Importance cho OLS và residual plots Phần 2
                        \item Tổng hợp nội dung báo cáo và chỉnh \texttt{README.md}
                    \end{itemize} & 100\% \\ \hline
\caption{Bảng đánh giá chi tiết}\\
\end{longtable}
\endgroup

\FloatBarrier


\section{Giới thiệu}


\subsection{Mục tiêu của báo cáo}
\begin{itemize}
    \item Trình bày nền tảng lý thuyết của bài toán data fitting và phương pháp OLS, bao gồm nghiệm OLS, ma trận Hat và các giả thiết Gauss--Markov.
    \item Cài đặt các thuật toán chính (OLS, Ridge, k-fold CV) và các kiểm định thống kê (F-test, t-test, khoảng tin cậy), đối chiếu với NumPy/sklearn để xác thực.
    \item Xây dựng pipeline tiền xử lý dữ liệu thực, đảm bảo đúng nguyên tắc fit/transform và tránh data leakage.
    \item So sánh hiệu suất nhiều mô hình trên tập test, phân tích phần dư và giải thích feature importance.
    \item Rút ra kết luận về mức độ phù hợp của mô hình tuyến tính trên dữ liệu Air Quality và định hướng cải thiện.
\end{itemize}


\subsection{Phạm vi trình bày}
\begin{itemize}
    \item Phần 1: lý thuyết OLS, minh họa bằng dữ liệu giả lập, kiểm định thống kê và phân tích phần dư.
    \item Phần 2: dữ liệu Air Quality, EDA, tiền xử lý, so sánh mô hình và đánh giá trên tập test.
    \item Kỹ thuật nâng cao: Bayesian Linear Regression để bổ sung điểm bonus.
\end{itemize}

\section{Phần 1: Lý thuyết Data Fitting và minh họa}

\subsection{Bài toán Data Fitting và mô hình hồi quy tuyến tính}
Bài toán data fitting tìm hàm $f:\mathbb{R}^p \to \mathbb{R}$ sao cho xấp xỉ tốt nhất mối quan hệ giữa $x_i$ và $y_i$ theo tiêu chí đã chọn. Trong hồi quy tuyến tính, ta giả thiết:
\begin{equation}
y_i = \beta_0 + \beta_1 x_{i1} + \cdots + \beta_p x_{ip} + \epsilon_i = x_i^T\beta + \epsilon_i
\end{equation}
Viết dưới dạng ma trận với $X\in\mathbb{R}^{n\times (p+1)}$ (cột đầu toàn 1):
\begin{equation}
y = X\beta + \epsilon
\end{equation}

\subsection{Nghiệm OLS và ước lượng phương sai}
OLS tối thiểu hóa tổng bình phương phần dư:
\begin{equation}
RSS(\beta)=\lVert y - X\beta \rVert_2^2 = \sum_{i=1}^{n} (y_i - x_i^T\beta)^2
\end{equation}
Nếu $X^T X$ khả nghịch, nghiệm OLS duy nhất là:
\begin{equation}
\hat{\beta}_{OLS}=(X^T X)^{-1}X^T y
\end{equation}
Ước lượng không chệch của phương sai nhiễu:
\begin{equation}
\hat{\sigma}^2 = \frac{RSS}{n-p-1}
\end{equation}

\subsection{Ma trận Hat và tính chất}
Ma trận chiếu (hat matrix) được định nghĩa:
\begin{equation}
H = X(X^T X)^{-1}X^T
\end{equation}
Các tính chất quan trọng: $H^2=H$ (lũy đẳng), $H^T=H$ (đối xứng), giá trị riêng chỉ là 0 hoặc 1, và $\text{rank}(H)=p+1$. Giá trị fitted là $\hat{y}=Hy$ và phần dư là $\hat{\epsilon}=(I-H)y$.

\subsection{Giả thiết Gauss--Markov và định lý BLUE}
Các giả thiết Gauss--Markov gồm: tuyến tính, không đa cộng tuyến hoàn hảo, ngoại sinh, đồng phương sai (và thường giả thiết chuẩn cho kiểm định). Khi thỏa GM1--GM4, ước lượng OLS là \textbf{BLUE} (Best Linear Unbiased Estimator):
\begin{equation}
\mathbb{E}[\hat{\beta}_{OLS}] = \beta, \quad \text{Var}(\hat{\beta}_{OLS}|X)=\sigma^2 (X^T X)^{-1}
\end{equation}

\subsection{Kiểm định thống kê và các chỉ số đánh giá}
Các chỉ số đánh giá và kiểm định quan trọng:
\begin{itemize}
    \item $RSS$, $TSS$, $R^2$ và $\overline{R}^2$ để đánh giá mức độ giải thích của mô hình.
    \item \textbf{F-test} kiểm định ý nghĩa toàn mô hình:
    \begin{equation}
    F = \frac{(TSS-RSS)/p}{RSS/(n-p-1)} \sim F_{p,n-p-1}
    \end{equation}
    \item \textbf{t-test} kiểm định ý nghĩa từng hệ số:
    \begin{equation}
    t_j = \frac{\hat{\beta}_j}{\hat{\sigma}\sqrt{[(X^T X)^{-1}]_{jj}}} \sim t_{n-p-1}
    \end{equation}
    \item \textbf{VIF} để phát hiện đa cộng tuyến:
    \begin{equation}
    VIF_j = \frac{1}{1-R_j^2}
    \end{equation}
\end{itemize}

\subsection{Minh họa định lý Gauss--Markov bằng Monte Carlo}
Nhóm thực hiện mô phỏng Monte Carlo với dữ liệu giả lập để kiểm chứng $\mathbb{E}[\hat{\beta}]=\beta$ và OLS có phương sai nhỏ nhất so với một ước lượng tuyến tính không chệch khác. Mỗi vòng lặp sinh nhiễu, fit OLS và một ước lượng thay thế (sử dụng tập con), sau đó so sánh kỳ vọng và phương sai thực nghiệm của hai ước lượng.

\begin{figure}[H]
    \centering
    \includegraphics[width=0.9\textwidth]{images/gauss_markov_mc.png}
    \caption{Minh họa Gauss--Markov bằng mô phỏng Monte Carlo.}
    \label{fig:gauss_markov_mc}
\end{figure}

\subsection{Phân tích phần dư cho dữ liệu giả lập}
Bốn biểu đồ chẩn đoán được sử dụng: Residuals vs Fitted, Q--Q Plot, Scale--Location, Cook's Distance. Kết quả cho thấy phần dư phân tán quanh 0, Q--Q gần tuyến tính, Scale--Location tương đối phẳng và không có điểm ảnh hưởng quá lớn, phù hợp với giả định OLS.

\begin{figure}[H]
    \centering
    \includegraphics[width=0.95\textwidth]{images/residual_plots_part1.png}
    \caption{Residual plots cho dữ liệu giả lập.}
    \label{fig:residual_part1}
\end{figure}

\section{Phần 2: Ứng dụng dữ liệu thực}

\subsection{Nguồn dữ liệu và điều kiện đề bài}

Bộ dữ liệu sử dụng là \textbf{Air Quality} (dữ liệu cảm biến chất lượng không khí theo thời gian). Dữ liệu bao gồm các biến cảm biến như \texttt{CO(GT)}, \texttt{PT08.S1(CO)}, \texttt{PT08.S2(NMHC)}, \texttt{NOx(GT)}, \texttt{T}, \texttt{RH}, \texttt{AH} cùng cột thời gian \texttt{Date/Time}. Target được chọn là \texttt{C6H6(GT)} (nồng độ benzene), là biến liên tục. Sau khi loại các dòng thiếu target, dữ liệu còn hơn 9000 quan sát và số biến đặc trưng lớn hơn 3.

Nguồn dữ liệu: \url{https://archive.ics.uci.edu/dataset/360/air+quality}.

Đối chiếu với yêu cầu đề bài:
\begin{itemize}
    \item Dữ liệu thực (real-world): thu thập từ cảm biến môi trường.
    \item Có missing values: nhiều cột có tỉ lệ missing vượt 5\%.
    \item Target liên tục: \texttt{C6H6(GT)}.
    \item Kích thước: $n \ge 200$, $p \ge 3$.
\end{itemize}
Như vậy bộ dữ liệu thỏa mãn đầy đủ điều kiện của đề bài.

\subsection{Tiền xử lý dữ liệu và DataPipeline}

\subsubsection{Phân tích missing values và outliers}
\subsubsection{Phân tích cơ chế giá trị khuyết thiếu (Missing Values)}
Dữ liệu từ các cảm biến đo lường thường xuyên bị thiếu sót. Qua phân tích, cơ chế missing value của bộ dữ liệu này được xác định chủ yếu là \textbf{MAR (Missing At Random - Khuyết thiếu ngẫu nhiên)}.

\textbf{Biện luận:} Sự cố của một cảm biến (ví dụ: cảm biến CO) thường không phụ thuộc vào chính nồng độ khí CO lúc đó, mà phụ thuộc vào các yếu tố ngoại cảnh khác như nhiệt độ hoặc độ ẩm quá cao/quá thấp khiến cảm biến bị lỗi tạm thời.

\textbf{Chiến lược xử lý:} Vì lý do trên, phương pháp \textbf{Listwise Deletion (xóa dòng)} bị loại bỏ vì sẽ làm mất mát một lượng lớn thông tin quan trọng. Thay vào đó, nhóm quyết định áp dụng chiến lược \textbf{Imputation (Điền khuyết)} để bảo toàn dữ liệu. Các cột có tỉ lệ missing quá cao ($> 50\%$) mới bị loại bỏ hoàn toàn do độ tin cậy không còn đảm bảo.

\subsubsection{Phân tích giá trị ngoại lai (Outliers)}
Biểu đồ Boxplot và phân phối cho thấy dữ liệu cảm biến chứa rất nhiều giá trị ngoại lai (những thời điểm nồng độ khí tăng vọt bất thường).

\textbf{Chiến lược xử lý:} Nhóm quyết định không xóa các điểm ngoại lai này (vì sẽ mâu thuẫn với nỗ lực điền khuyết bảo toàn dữ liệu ở trên), mà sử dụng phương pháp \textbf{Winsorization (Cắt ngọn - IQR Capping)}. Ngoài ra, việc tồn tại nhiều outliers cũng là cơ sở để nhóm quyết định dùng \textbf{Median (Trung vị)} thay vì Mean (Trung bình) trong quá trình điền khuyết các biến số thực, do Median mang tính bền vững (robust) hơn trước outlier.

Công thức IQR capping:
\begin{align*}
IQR &= Q_3 - Q_1 \\
\text{Lower} &= Q_1 - 1.5\times IQR \\
\text{Upper} &= Q_3 + 1.5\times IQR
\end{align*}

\subsubsection{Quy trình DataPipeline}
Pipeline gồm 6 bước, thực hiện theo đúng thứ tự và chỉ \textbf{fit trên train}:
\begin{enumerate}
    \item \textbf{Feature engineering:} trích xuất \texttt{Month}, \texttt{DayOfWeek}, \texttt{Hour} từ \texttt{Date/Time}, sau đó loại bỏ cột gốc.
    \item \textbf{Loại cột thiếu quá nhiều:} các cột có tỉ lệ missing vượt ngưỡng 50\% được loại bỏ.
    \item \textbf{Imputation:} biến số thực điền bằng median (robust với outliers), biến phân loại điền bằng mode. Với biến thời gian, dùng median để đảm bảo giá trị hợp lệ.
    \item \textbf{Outlier handling:} áp dụng IQR capping cho các biến số (không áp dụng cho biến thời gian).
    \item \textbf{One-hot encoding:} mã hóa các biến phân loại, đồng thời reindex để test có cùng cấu trúc cột như train.
    \item \textbf{Chuẩn hóa Z-score:} $x_{std} = (x-\mu)/\sigma$ để đưa các biến về cùng thang đo.
\end{enumerate}

Nguyên tắc chống data leakage:
\begin{itemize}
    \item Tất cả tham số (median/mode, IQR bounds, danh sách cột, mean/std) chỉ học từ train.
    \item Test chỉ được transform bằng tham số đã học.
\end{itemize}

\subsection{So sánh mô hình và đánh giá}

\subsection{Ridge Regression và Ridge Trace}
OLS tìm nghiệm bằng cách tối thiểu hóa tổng bình phương phần dư:
\begin{equation}
\hat{\beta}_{OLS} = (X^T X)^{-1}X^T y
\end{equation}

Khi các cột trong $X$ tương quan mạnh, $X^T X$ có thể gần suy biến khiến hệ số OLS kém ổn định. Ridge Regression thêm thành phần phạt $L_2$:
\begin{equation}
\hat{\beta}_{ridge} = \arg\min_{\beta}\left\{\lVert y - X\beta \rVert_2^2 + \lambda \lVert \beta \rVert_2^2\right\}
\end{equation}
Nghiệm đóng của Ridge là:
\begin{equation}
\hat{\beta}_{ridge} = (X^T X + \lambda I')^{-1} X^T y
\end{equation}
trong đó $I'$ là ma trận đơn vị nhưng phần tử $I'_{00}=0$ để không regularize hệ số chặn.

Ridge Trace là biểu đồ thể hiện sự thay đổi của từng hệ số theo $\log_{10}(\lambda)$. Khi $\lambda$ tăng, các hệ số bị co dần về 0; hệ số ổn định qua nhiều giá trị $\lambda$ thường có vai trò quan trọng hơn. Ridge Trace giúp quan sát trực quan mức độ co hệ số và hỗ trợ lựa chọn vùng $\lambda$ phù hợp, tránh underfit.

\subsection{k-fold Cross-Validation}
k-fold CV chia tập train thành $k$ phần gần bằng nhau. Mỗi lần dùng $k-1$ phần để huấn luyện và 1 phần để kiểm tra, lặp $k$ lần. Điểm CV-MSE được tính:
\begin{equation}
CV_{(k)} = \frac{1}{k}\sum_{i=1}^{k} MSE_i
\end{equation}
Trong dự án, k-fold CV được dùng để chọn $\lambda$ tối ưu cho Ridge. Mỗi giá trị $\lambda$ được đánh giá bằng CV-MSE, và $\lambda$ có CV-MSE nhỏ nhất được chọn.

Kết quả được lấy từ script \texttt{part2/model\_comparison.py} với seed 42, chia train/test 80/20 và chọn lambda Ridge bằng 5-fold CV trên tập train. Các chỉ số đánh giá:
\begin{align*}
MAE &= \frac{1}{n}\sum_i |y_i-\hat{y}_i| \\
RMSE &= \sqrt{\frac{1}{n}\sum_i (y_i-\hat{y}_i)^2} \\
R^2 &= 1 - \frac{RSS}{TSS}
\end{align*}

\begin{table}[H]
\centering
\begin{tabular}{lrrr}
\hline
Model & MAE & RMSE & $R^2$ \\
\hline
OLS & 0.889981 & 1.446089 & 0.963608 \\
OLS (Variable Selection) & 1.428491 & 2.072625 & 0.925242 \\
Ridge ($\lambda=0.1$) & 0.889957 & 1.446081 & 0.963608 \\
Bayesian Linear Regression & 0.889976 & 1.446087 & 0.963608 \\
\hline
\end{tabular}
\caption{So sánh hiệu suất các mô hình trên tập test Air Quality.}
\end{table}

\subsection{Nhận xét}
\subsubsection{OLS cơ bản}
OLS cơ bản sử dụng toàn bộ biến sau tiền xử lý. Mô hình đạt MAE = 0.889981, RMSE = 1.446089 và $R^2 = 0.963608$ trên tập test. Điều này cho thấy quan hệ tuyến tính giữa các tín hiệu cảm biến và nồng độ benzene là rất mạnh. Tuy nhiên, OLS nhạy với đa cộng tuyến nên hệ số có thể kém ổn định khi các cảm biến tương quan cao.

\subsubsection{OLS chọn biến}
OLS chọn biến loại bỏ các biến có VIF > 10 hoặc p-value > 0.05. Trong lần chạy này, các biến \texttt{PT08.S2(NMHC)} và \texttt{T} bị loại do VIF cao. Hiệu suất giảm đáng kể (MAE = 1.428491, RMSE = 2.072625, $R^2 = 0.925242$), cho thấy việc loại biến làm mất một phần thông tin dự báo quan trọng, đặc biệt khi các biến đa cộng tuyến vẫn mang tín hiệu mạnh với target.

\begin{table}[H]
\centering
\begin{tabular}{l l l r}
\hline
Bước & Biến bị loại & Lý do & Giá trị \\
\hline
1 & PT08.S2(NMHC) & VIF & 31.742536 \\
2 & T & VIF & 14.549521 \\
\hline
\end{tabular}
\caption{Các biến bị loại trong OLS chọn biến.}
\end{table}

\subsubsection{Ridge Regression}
Ridge được chọn với $\lambda=0.1$ bằng k-fold CV. Mô hình đạt MAE = 0.889957 và RMSE = 1.446081 (tốt nhất nhưng chỉ nhỉnh hơn OLS rất nhỏ). Ridge giữ lại toàn bộ biến và co nhỏ hệ số, nên phù hợp khi muốn giảm variance mà không loại bỏ tín hiệu.

\subsubsection{Bayesian Linear Regression}
Bayesian Linear Regression đạt MAE = 0.889976, RMSE = 1.446087 và $R^2 = 0.963608$. Kết quả gần OLS/Ridge vì prior tương đối yếu sau khi chuẩn hóa dữ liệu. Ưu điểm là mô hình cung cấp posterior covariance, giúp diễn giải độ bất định của hệ số.

\subsubsection{Tổng hợp}
OLS, Ridge và Bayesian Linear Regression có hiệu suất gần như tương đương, phản ánh dữ liệu phù hợp với mô hình tuyến tính. OLS chọn biến giúp mô hình gọn hơn nhưng giảm độ chính xác. Ridge là lựa chọn cân bằng nếu ưu tiên tính ổn định của hệ số.

\subsection{Feature importance và residual plots}
Hệ số hồi quy sau chuẩn hóa được dùng để biểu diễn \textbf{feature importance}. Các biến có |hệ số| lớn nhất là nhóm ảnh hưởng mạnh nhất tới \texttt{C6H6(GT)}; trong dữ liệu này, các cảm biến liên quan đến hydrocarbon như \texttt{PT08.S2(NMHC)} và \texttt{CO(GT)} thường có tác động mạnh, bên cạnh các biến điều kiện môi trường.

\begin{figure}[H]
    \centering
    \includegraphics[width=0.9\textwidth]{images/feature_importance_ols.png}
    \caption{Feature importance theo hệ số OLS đã chuẩn hóa.}
    \label{fig:feature_importance_ols}
\end{figure}

Đồng thời, 4 biểu đồ phần dư (Residuals vs Fitted, Q--Q, Scale--Location, Cook's Distance) được dùng để kiểm tra giả thiết tuyến tính, chuẩn và đồng phương sai. Kết quả cho thấy phần dư phân tán quanh 0 và không có điểm ảnh hưởng quá lớn, phù hợp với giả định OLS.

\begin{figure}[H]
    \centering
    \includegraphics[width=0.95\textwidth]{images/residual_plots_part2.png}
    \caption{Residual plots cho dữ liệu Air Quality.}
    \label{fig:residual_part2}
\end{figure}

\subsection{Kỹ thuật nâng cao: Bayesian Linear Regression}

Bayesian Linear Regression được dùng để lấy điểm bonus. Mô hình giữ cấu trúc tuyến tính:
\begin{align*}
y &= X\beta + \epsilon, \quad \epsilon \sim \mathcal{N}(0, \sigma^2 I) \\
\beta &\sim \mathcal{N}(m_0, S_0)
\end{align*}
Sau khi quan sát dữ liệu, posterior của hệ số là:
\begin{align*}
S_n &= \left(S_0^{-1} + \frac{1}{\sigma^2}X^T X\right)^{-1} \\
m_n &= S_n\left(S_0^{-1}m_0 + \frac{1}{\sigma^2}X^T y\right)
\end{align*}
Trong triển khai hiện tại, prior mean là vector 0, prior variance = 100 cho các hệ số và intercept prior variance = $10^{12}$ để gần như không regularize intercept. Khi prior yếu, posterior mean gần OLS; khi prior mạnh, hệ số bị shrink. Lựa chọn Bayesian Linear Regression phù hợp với đề bài vì vẫn giữ cấu trúc tuyến tính, dễ giải thích hệ số và không cần ma trận Gram lớn như Kernel Regression.

\subsection{Lý do chọn Bayesian Linear Regression thay vì Kernel Regression}
Kernel Regression có khả năng mô hình hóa quan hệ phi tuyến nhưng đòi hỏi xây dựng và nghịch đảo ma trận Gram kích thước $n \times n$, không phù hợp khi $n$ lớn như bộ Air Quality. Trong khi đó, Bayesian Linear Regression chỉ cần nghịch đảo ma trận kích thước $(p+1)\times(p+1)$ nên phù hợp với helper ma trận tự cài. Ngoài ra, dữ liệu Air Quality cho thấy quan hệ tuyến tính rất mạnh với target (OLS/Ridge đạt $R^2$ khoảng 0.964), nên lợi ích kỳ vọng từ Kernel Regression không lớn so với chi phí triển khai. Bayesian Linear Regression cũng giữ được hệ số tuyến tính, dễ giải thích feature importance theo yêu cầu đề bài.

\section{Phân tích các hàm chính}

Mục này trình bày theo một khung thống nhất để người đọc dễ theo dõi: bối cảnh sử dụng, cơ sở lý thuyết, đầu vào, đầu ra, quy trình xử lý, và kiểm thử. Mọi hàm đều dùng biểu diễn ma trận bằng danh sách hai chiều để đồng nhất với các helper ma trận tự cài đặt.

\subsection{Nhóm hàm OLS và kiểm định thống kê}

\subsubsection{\code{ols\_fit}}
\begin{itemize}[leftmargin=*]
    \item \textbf{Bối cảnh sử dụng:} Hàm lõi cho toàn bộ phần OLS, được tái sử dụng trong kiểm định thống kê, VIF và các phần lựa chọn biến.
    \item \textbf{Cơ sở lý thuyết:} Nghiệm đóng của OLS và ước lượng phương sai nhiễu $\hat{\sigma}^2 = RSS/(n-p-1)$.
    \item \textbf{Đầu vào:} $X$ là ma trận $n\times(p+1)$ có cột chặn, $y$ là vector cột $n\times1$.
    \item \textbf{Đầu ra:} $\hat{\beta}$ dạng $(p+1)\times1$ và $\hat{\sigma}^2$.
    \item \textbf{Quy trình xử lý:} Kiểm tra hợp lệ của ma trận, tính $X^T X$ và nghịch đảo, suy ra $\hat{\beta}$, sau đó tính $\hat{y}$, phần dư và $RSS$ để suy ra $\hat{\sigma}^2$.
    \item \textbf{Lưu ý:} Nếu bậc tự do $n-(p+1)$ không dương, hàm báo lỗi để tránh ước lượng sai.
\end{itemize}
\begin{figure}[H]
    \centering
    \includegraphics[width=0.9\textwidth]{images/test_ols_fit.png}
    \caption{Ảnh kiểm thử hàm \texttt{ols\_fit}.}
\end{figure}

\subsubsection{\code{hat\_matrix}}
\begin{itemize}[leftmargin=*]
    \item \textbf{Bối cảnh sử dụng:} Phục vụ phân tích phần dư, tính đòn bẩy và khoảng Cook.
    \item \textbf{Cơ sở lý thuyết:} Ma trận chiếu $H=X(X^T X)^{-1}X^T$ với tính chất đối xứng và lũy đẳng.
    \item \textbf{Đầu vào:} Ma trận $X$ đã có cột chặn.
    \item \textbf{Đầu ra:} Ma trận $H$ kích thước $n\times n$.
    \item \textbf{Quy trình xử lý:} Tính $X^T X$, nghịch đảo và nhân lần lượt để tạo $H$.
    \item \textbf{Lưu ý:} Nếu $X^T X$ suy biến thì phép nghịch đảo sẽ lỗi, vì vậy dữ liệu phải không đa cộng tuyến hoàn hảo.
\end{itemize}
\begin{figure}[H]
    \centering
    \includegraphics[width=0.9\textwidth]{images/test_hat_matrix.png}
    \caption{Ảnh kiểm thử hàm \texttt{hat\_matrix}.}
\end{figure}

\subsubsection{\code{model\_metrics}}
\begin{itemize}[leftmargin=*]
    \item \textbf{Bối cảnh sử dụng:} Tổng hợp các chỉ số đánh giá và kiểm định cho mô hình OLS.
    \item \textbf{Cơ sở lý thuyết:} $R^2$, $\overline{R}^2$ và kiểm định F dựa trên $RSS$ và $TSS$.
    \item \textbf{Đầu vào:} $y$, $\hat{y}$ và số biến giải thích $p$ không tính cột chặn.
    \item \textbf{Đầu ra:} Từ điển gồm $RSS$, $TSS$, $R^2$, $\overline{R}^2$, thống kê F và giá trị p.
    \item \textbf{Quy trình xử lý:} Tính trung bình $y$, suy ra $RSS$ và $TSS$, sau đó tính các chỉ số và kiểm định F.
    \item \textbf{Lưu ý:} Nếu mẫu quá nhỏ làm mẫu số gần 0, hàm trả về F vô hạn và p bằng 0 để tránh lỗi số học.
\end{itemize}
\begin{figure}[H]
    \centering
    \includegraphics[width=0.9\textwidth]{images/test_model_metrics.png}
    \caption{Ảnh kiểm thử hàm \texttt{model\_metrics}.}
\end{figure}

\subsubsection{\code{coef\_inference}}
\begin{itemize}[leftmargin=*]
    \item \textbf{Bối cảnh sử dụng:} Suy luận thống kê cho từng hệ số để phục vụ chọn biến và diễn giải mô hình.
    \item \textbf{Cơ sở lý thuyết:} Ma trận hiệp phương sai của $\hat{\beta}$ và kiểm định t theo phân phối Student.
    \item \textbf{Đầu vào:} $X$, $y$, $\hat{\beta}$ và $\hat{\sigma}^2$ từ \texttt{ols\_fit}.
    \item \textbf{Đầu ra:} Danh sách sai số chuẩn, thống kê t, giá trị p và khoảng tin cậy 95\%.
    \item \textbf{Quy trình xử lý:} Tính $(X^T X)^{-1}$, suy ra sai số chuẩn, thống kê t, sau đó lấy ngưỡng tới hạn 0.975 để lập khoảng tin cậy hai phía.
    \item \textbf{Lưu ý:} Diễn giải thống kê t dựa trên giả định phần dư có phân phối gần chuẩn.
\end{itemize}
\begin{figure}[H]
    \centering
    \includegraphics[width=0.9\textwidth]{images/test_coef_inference.png}
    \caption{Ảnh kiểm thử hàm \texttt{coef\_inference}.}
\end{figure}

\subsubsection{\code{vif}}
\begin{itemize}[leftmargin=*]
    \item \textbf{Bối cảnh sử dụng:} Đo đa cộng tuyến để hỗ trợ loại biến không ổn định.
    \item \textbf{Cơ sở lý thuyết:} $VIF_j=1/(1-R_j^2)$ với $R_j^2$ từ hồi quy biến $j$ theo các biến còn lại.
    \item \textbf{Đầu vào:} Ma trận $X$ có cột chặn.
    \item \textbf{Đầu ra:} Danh sách VIF theo từng biến độc lập (bỏ qua cột chặn).
    \item \textbf{Quy trình xử lý:} Với mỗi biến, tạo bài toán hồi quy phụ, tính $R^2$ và suy ra VIF.
    \item \textbf{Lưu ý:} Nếu $1-R^2$ rất nhỏ thì VIF tiến tới vô hạn, phản ánh đa cộng tuyến mạnh.
\end{itemize}
\begin{figure}[H]
    \centering
    \includegraphics[width=0.9\textwidth]{images/test_vif.png}
    \caption{Ảnh kiểm thử hàm \texttt{vif}.}
\end{figure}

\subsection{Nhóm hàm Ridge và k-fold CV}

\subsubsection{\code{ridge\_fit}}
\begin{itemize}[leftmargin=*]
    \item \textbf{Bối cảnh sử dụng:} Tính nghiệm Ridge để giảm phương sai khi các biến có tương quan cao.
    \item \textbf{Cơ sở lý thuyết:} $\hat{\beta}=(X^T X+\lambda I')^{-1}X^T y$, trong đó $I'$ bỏ phạt hệ số chặn.
    \item \textbf{Đầu vào:} $X$, $y$ và tham số phạt $\lambda$ không âm.
    \item \textbf{Đầu ra:} $\hat{\beta}$ Ridge dạng $(p+1)\times1$.
    \item \textbf{Quy trình xử lý:} Tạo ma trận đơn vị, đặt phần tử chặn bằng 0, cộng với $X^T X$ rồi nghịch đảo để suy ra nghiệm.
    \item \textbf{Lưu ý:} Khi $\lambda=0$, nghiệm Ridge trùng OLS.
\end{itemize}
\begin{figure}[H]
    \centering
    \includegraphics[width=0.9\textwidth]{images/test_ridge_fit.png}
    \caption{Ảnh kiểm thử hàm \texttt{ridge\_fit}.}
\end{figure}

\subsubsection{\codet{plot\_ridge\_trace}}
\begin{itemize}[leftmargin=*]
    \item \textbf{Bối cảnh sử dụng:} Minh họa sự co rút hệ số khi tăng $\lambda$ để chọn vùng ổn định.
    \item \textbf{Đầu vào:} $X$, $y$, danh sách $\lambda$; tùy chọn tên biến và số biến tối đa để hiển thị.
    \item \textbf{Đầu ra:} Danh sách quỹ đạo hệ số và biểu đồ đường theo $\log_{10}(\lambda)$.
    \item \textbf{Quy trình xử lý:} Với mỗi $\lambda$, gọi \texttt{ridge\_fit}, gom hệ số, sau đó vẽ từng đường hệ số.
    \item \textbf{Lưu ý:} Nếu không truyền $\lambda$ thì hàm sinh 100 giá trị theo thang lũy thừa từ $10^{-4}$ đến $10^4$.
\end{itemize}
\begin{figure}[H]
    \centering
    \includegraphics[width=0.9\textwidth]{images/test_ridge_trace.png}
    \caption{Ảnh kiểm thử hàm \texttt{plot\_ridge\_trace}.}
\end{figure}

\subsubsection{\code{kfold\_cv}}
\begin{itemize}[leftmargin=*]
    \item \textbf{Bối cảnh sử dụng:} Đánh giá độ ổn định của mô hình trên nhiều cách chia dữ liệu.
    \item \textbf{Đầu vào:} $X$, $y$, số $k$ và hàm huấn luyện tùy chọn trả về $\hat{\beta}$.
    \item \textbf{Đầu ra:} MSE trung bình và danh sách MSE theo từng phần.
    \item \textbf{Quy trình xử lý:} Trộn chỉ số theo seed cố định, chia k phần, lặp qua từng phần làm kiểm tra, tính MSE và lấy trung bình.
    \item \textbf{Lưu ý:} Giá trị $k$ phải trong khoảng $[2,n]$ để đảm bảo mỗi phần có dữ liệu.
\end{itemize}
\begin{figure}[H]
    \centering
    \includegraphics[width=0.9\textwidth]{images/test_kfold_cv.png}
    \caption{Ảnh kiểm thử hàm \texttt{kfold\_cv}.}
\end{figure}

\subsection{Nhóm hàm phân tích phần dư}

\subsubsection{\code{residual\_plots}}
\begin{itemize}[leftmargin=*]
    \item \textbf{Bối cảnh sử dụng:} Kiểm tra giả định tuyến tính, chuẩn và đồng phương sai cho mô hình.
    \item \textbf{Cơ sở lý thuyết:} Dựa trên phần dư chuẩn hóa, đòn bẩy từ ma trận chiếu và khoảng Cook.
    \item \textbf{Đầu vào:} $X$, $y$, $\hat{\beta}$.
    \item \textbf{Đầu ra:} Một bộ hình gồm bốn biểu đồ chẩn đoán và cặp đối tượng hình, trục để tái sử dụng.
    \item \textbf{Quy trình xử lý:} Tính $\hat{y}$, phần dư, $\hat{\sigma}^2$, đòn bẩy, phần dư chuẩn hóa và khoảng Cook, sau đó vẽ bốn biểu đồ.
    \item \textbf{Lưu ý:} Ngưỡng khoảng Cook dùng $4/n$ để gợi ý điểm ảnh hưởng lớn.
\end{itemize}
\begin{figure}[H]
    \centering
    \includegraphics[width=0.95\textwidth]{images/test_residual_plots.png}
    \caption{Ảnh kiểm thử hàm \texttt{residual\_plots}.}
\end{figure}

\subsection{Nhóm hàm tiền xử lý dữ liệu}

\subsubsection{\code{DataPipeline.fit}}
\begin{itemize}[leftmargin=*]
    \item \textbf{Bối cảnh sử dụng:} Học toàn bộ tham số tiền xử lý từ tập huấn luyện để tránh rò rỉ dữ liệu.
    \item \textbf{Đầu vào:} Bảng dữ liệu thô chưa tách target.
    \item \textbf{Đầu ra:} Đối tượng pipeline đã lưu cột bị loại, giá trị điền khuyết, biên cắt IQR, cấu trúc cột mã hóa và tham số chuẩn hóa.
    \item \textbf{Quy trình xử lý:} Trích xuất tháng, thứ trong tuần, giờ; loại cột thiếu quá ngưỡng; điền khuyết theo chiến lược; cắt ngọn outlier bằng IQR; mã hóa một nóng; lưu trung bình và độ lệch chuẩn cho chuẩn hóa.
    \item \textbf{Lưu ý:} Biến thời gian không bị cắt ngọn để giữ cấu trúc chu kỳ và dùng giá trị trung vị để tránh ảnh hưởng của nhiễu.
\end{itemize}

\subsubsection{\code{DataPipeline.transform}}
\begin{itemize}[leftmargin=*]
    \item \textbf{Bối cảnh sử dụng:} Áp dụng lại đúng các tham số đã học trên tập mới.
    \item \textbf{Đầu vào:} Bảng dữ liệu cần xử lý.
    \item \textbf{Đầu ra:} Bảng dữ liệu đã được biến đổi và chuẩn hóa theo cấu trúc của tập huấn luyện.
    \item \textbf{Quy trình xử lý:} Trích xuất thời gian, bỏ cột đã loại, điền khuyết theo giá trị đã học, cắt ngọn theo biên IQR, mã hóa một nóng, sau đó chuẩn hóa theo Z-score.
    \item \textbf{Lưu ý:} Có bước tái lập cột theo cấu trúc đã học để đảm bảo tập kiểm tra không thiếu cột.
\end{itemize}
\begin{figure}[H]
    \centering
    \includegraphics[width=0.9\textwidth]{images/test_data_pipeline.png}
    \caption{Ảnh kiểm thử lớp \texttt{DataPipeline}.}
\end{figure}

\subsection{Nhóm hàm Bayesian Linear Regression}

\subsubsection{\code{bayesian\_linear\_fit}}
\begin{itemize}[leftmargin=*]
    \item \textbf{Bối cảnh sử dụng:} Bổ sung mô hình Bayes với prior Gaussian, vẫn giữ cấu trúc tuyến tính.
    \item \textbf{Cơ sở lý thuyết:} Hậu nghiệm Gaussian với $S_n=(S_0^{-1}+X^T X/\sigma^2)^{-1}$ và $m_n=S_n(S_0^{-1}m_0+X^T y/\sigma^2)$.
    \item \textbf{Đầu vào:} $X$, $y$, tham số prior, và $\sigma^2$ có thể ước lượng từ OLS.
    \item \textbf{Đầu ra:} Posterior mean, posterior covariance, cùng các tham số prior và $\sigma^2$.
    \item \textbf{Quy trình xử lý:} Kiểm tra kích thước, tạo prior precision theo phương sai, cộng với precision từ dữ liệu, nghịch đảo để lấy hậu nghiệm.
    \item \textbf{Lưu ý:} Intercept có phương sai prior rất lớn để gần như không bị điều chuẩn.
\end{itemize}

\subsubsection{\code{bayesian\_predict}}
\begin{itemize}[leftmargin=*]
    \item \textbf{Bối cảnh sử dụng:} Dự đoán theo kỳ vọng hậu nghiệm.
    \item \textbf{Đầu vào:} $X$ và từ điển posterior từ \texttt{bayesian\_linear\_fit}.
    \item \textbf{Đầu ra:} Vector dự đoán $\hat{y}$ dạng $n\times1$.
    \item \textbf{Quy trình xử lý:} Nhân ma trận $X$ với $\mathbb{E}[\beta|X,y]$.
\end{itemize}

\subsubsection{\code{credible\_intervals}}
\begin{itemize}[leftmargin=*]
    \item \textbf{Bối cảnh sử dụng:} Ước lượng khoảng tin cậy của hệ số theo hậu nghiệm.
    \item \textbf{Đầu vào:} Posterior và mức tin cậy $\alpha$.
    \item \textbf{Đầu ra:} Danh sách khoảng tin cậy cho từng hệ số.
    \item \textbf{Quy trình xử lý:} Lấy phương sai trên đường chéo của ma trận hậu nghiệm, dùng hệ số z để tạo khoảng đối xứng.
    \item \textbf{Lưu ý:} Mặc định mức 95\% dùng hệ số 1.96 để báo cáo gọn.
\end{itemize}
\begin{figure}[H]
    \centering
    \includegraphics[width=0.9\textwidth]{images/test_bayesian_linear.png}
    \caption{Ảnh kiểm thử nhóm hàm Bayesian Linear Regression.}
\end{figure}

\subsection{Nhóm hàm so sánh mô hình}

\subsubsection{\code{train\_and\_compare}}
\begin{itemize}[leftmargin=*]
    \item \textbf{Bối cảnh sử dụng:} Quy trình tổng hợp từ tiền xử lý đến huấn luyện và so sánh mô hình.
    \item \textbf{Đầu vào:} Đường dẫn dữ liệu, số k cho chia k phần, danh sách $\lambda$ tùy chọn và cờ vẽ đồ thị.
    \item \textbf{Đầu ra:} Bảng chỉ số trên tập kiểm tra, kết quả chọn $\lambda$, mô hình đã huấn luyện và dự đoán của từng mô hình.
    \item \textbf{Quy trình xử lý:} Đọc dữ liệu, tách target, chia tập, chạy pipeline, thêm cột chặn, chọn $\lambda$ bằng CV, huấn luyện OLS, chọn biến, Ridge và Bayes, sau đó đánh giá MAE, RMSE, $R^2$.
    \item \textbf{Lưu ý:} Có thể bật vẽ đồ thị để quan sát đường chọn $\lambda$ và quỹ đạo hệ số.
\end{itemize}
\begin{figure}[H]
    \centering
    \includegraphics[width=0.9\textwidth]{images/test_model_comparison.png}
    \caption{Ảnh kiểm thử hàm \texttt{train\_and\_compare}.}
\end{figure}

\section{Kết luận}

Kết quả thực nghiệm cho thấy mô hình tuyến tính mô tả tốt dữ liệu Air Quality với $R^2 \approx 0.964$. Ridge đạt MAE/RMSE tốt nhất nhưng chênh lệch rất nhỏ so với OLS và Bayesian. OLS chọn biến giúp mô hình gọn hơn nhưng làm giảm hiệu suất đáng kể do loại bỏ các biến có tín hiệu dự báo mạnh.

Các điểm đạt được:
\begin{itemize}
    \item Cài đặt đầy đủ các thuật toán yêu cầu từ công thức toán học và có kiểm chứng.
    \item Pipeline tiền xử lý đảm bảo fit/transform đúng chuẩn, tránh data leakage.
    \item So sánh nhiều mô hình và phân tích phần dư đáp ứng tiêu chí đề bài.
\end{itemize}

Hạn chế và hướng mở rộng:
\begin{itemize}
    \item Có thể thử thêm biến tương tác hoặc biến phi tuyến để kiểm tra cải thiện hiệu suất.
    \item Đánh giá độ ổn định của chọn biến qua nhiều lần split hoặc k-fold trên toàn bộ pipeline.
    \item Kết hợp thêm các kỹ thuật regularization khác nếu mục tiêu ưu tiên dự báo.
\end{itemize}

\section*{Tài liệu tham khảo}
\addcontentsline{toc}{section}{Tài liệu tham khảo}
\begin{enumerate}
    \item Gilbert Strang. \textit{Introduction to Linear Algebra}, 6th ed. Wellesley-Cambridge Press, 2023.
    \item Gareth James, Daniela Witten, Trevor Hastie, Robert Tibshirani. \textit{An Introduction to Statistical Learning}, 2nd ed. Springer, 2021. \url{https://www.statlearning.com}
    \item Trevor Hastie, Robert Tibshirani, Jerome Friedman. \textit{The Elements of Statistical Learning}, 2nd ed. Springer, 2009. \url{https://hastie.su.domains/ElemStatLearn/}
    \item Christopher M. Bishop. \textit{Pattern Recognition and Machine Learning}. Springer, 2006.
    \item Kevin P. Murphy. \textit{Probabilistic Machine Learning: An Introduction}. MIT Press, 2022. \url{https://probml.github.io/pml-book/book1.html}
    \item Jake VanderPlas. \textit{Python Data Science Handbook}. O'Reilly, 2016. \url{https://jakevdp.github.io/PythonDataScienceHandbook/}
    \item Wes McKinney. \textit{Python for Data Analysis}, 3rd ed. O'Reilly, 2022.
\end{enumerate}

\subsection*{Phụ lục A: Bảng số liệu và kết quả bổ sung}
\addcontentsline{toc}{subsection}{Phụ lục A: Bảng số liệu và kết quả bổ sung}

\subsubsection*{A1. Thống kê mô tả các biến số}
\begin{table}[H]
    \centering
    \scriptsize
    \begin{tabular}{lrrrrrrrr}
    \toprule
     & count & mean & std & min & 25\% & 50\% & 75\% & max \\
    \midrule
    CO(GT) & 7344.000000 & 2.129700 & 1.436500 & 0.100000 & 1.100000 & 1.800000 & 2.800000 & 11.900000 \\
    PT08.S1(CO) & 8991.000000 & 1099.833200 & 217.080000 & 647.000000 & 937.000000 & 1063.000000 & 1231.000000 & 2040.000000 \\
    NMHC(GT) & 887.000000 & 218.607700 & 206.615100 & 7.000000 & 66.000000 & 145.000000 & 297.000000 & 1189.000000 \\
    PT08.S2(NMHC) & 8991.000000 & 939.153400 & 266.831400 & 383.000000 & 734.500000 & 909.000000 & 1116.000000 & 2214.000000 \\
    NOx(GT) & 7396.000000 & 242.189300 & 206.312000 & 2.000000 & 97.000000 & 178.000000 & 321.000000 & 1479.000000 \\
    PT08.S3(NOx) & 8991.000000 & 835.493600 & 256.817300 & 322.000000 & 658.000000 & 806.000000 & 969.500000 & 2683.000000 \\
    NO2(GT) & 7393.000000 & 112.145100 & 47.629100 & 2.000000 & 77.000000 & 109.000000 & 140.000000 & 333.000000 \\
    PT08.S4(NO2) & 8991.000000 & 1456.264600 & 346.206800 & 551.000000 & 1227.000000 & 1463.000000 & 1674.000000 & 2775.000000 \\
    PT08.S5(O3) & 8991.000000 & 1022.906100 & 398.484300 & 221.000000 & 731.500000 & 963.000000 & 1273.5000 00 & 2523.000000 \\
    T & 8991.000000 & 18.317800 & 8.832100 & -1.900000 & 11.800000 & 17.800000 & 24.400000 & 44.600000 \\
    RH & 8991.000000 & 49.234200 & 17.316900 & 9.200000 & 35.800000 & 49.600000 & 62.500000 & 88.700000 \\
    AH & 8991.000000 & 1.025500 & 0.403800 & 0.184700 & 0.736800 & 0.995400 & 1.313700 & 2.231000 \\
    C6H6(GT) & 8991.000000 & 10.083100 & 7.449800 & 0.100000 & 4.400000 & 8.200000 & 14.000000 & 63.700000 \\
    \bottomrule
    \end{tabular}
    \caption{Thống kê mô tả các biến số trong Air Quality (sau khi loại dòng thiếu target).}
\end{table}

\subsubsection*{A2. Ma trận tương quan (các biến quan trọng và target)}
\begin{table}[H]
    \centering
    \small
    \resizebox{0.95\textwidth}{!}{
    \begin{tabular}{lrrrrrrrrrrr}
    \toprule
     & PT08.S2(NMHC) & PT08.S3(NOx) & T & AH & PT08.S1(CO) & PT08.S4(NO2) & RH & Hour & NOx(GT) & PT08.S5(O3) & C6H6(GT) \\
    \midrule
    PT08.S2(NMHC) & 1.000000 & -0.826000 & 0.239000 & 0.185000 & 0.894000 & 0.773000 & -0.086000 & 0.397000 & 0.644000 & 0.881000 & 0.976000 \\
    PT08.S3(NOx) & -0.826000 & 1.000000 & -0.136000 & -0.229000 & -0.807000 & -0.550000 & -0.066000 & -0.338000 & -0.675000 & -0.827000 & -0.766000 \\
    T & 0.239000 & -0.136000 & 1.000000 & 0.658000 & 0.046000 & 0.566000 & -0.579000 & 0.204000 & -0.267000 & -0.031000 & 0.193000 \\
    AH & 0.185000 & -0.229000 & 0.658000 & 1.000000 & 0.132000 & 0.634000 & 0.166000 & -0.023000 & -0.161000 & 0.071000 & 0.167000 \\
    PT08.S1(CO) & 0.894000 & -0.807000 & 0.046000 & 0.132000 & 1.000000 & 0.674000 & 0.118000 & 0.335000 & 0.666000 & 0.902000 & 0.879000 \\
    PT08.S4(NO2) & 0.773000 & -0.550000 & 0.566000 & 0.634000 & 0.674000 & 1.000000 & -0.031000 & 0.223000 & 0.188000 & 0.589000 & 0.757000 \\
    RH & -0.086000 & -0.066000 & -0.579000 & 0.166000 & 0.118000 & -0.031000 & 1.000000 & -0.296000 & 0.198000 & 0.134000 & -0.053000 \\
    Hour & 0.397000 & -0.338000 & 0.204000 & -0.023000 & 0.335000 & 0.223000 & -0.296000 & 1.000000 & 0.244000 & 0.250000 & 0.351000 \\
    NOx(GT) & 0.644000 & -0.675000 & -0.267000 & -0.161000 & 0.666000 & 0.188000 & 0.198000 & 0.244000 & 1.000000 & 0.741000 & 0.632000 \\
    PT08.S5(O3) & 0.881000 & -0.827000 & -0.031000 & 0.071000 & 0.902000 & 0.589000 & 0.134000 & 0.250000 & 0.741000 & 1.000000 & 0.862000 \\
    C6H6(GT) & 0.976000 & -0.766000 & 0.193000 & 0.167000 & 0.879000 & 0.757000 & -0.053000 & 0.351000 & 0.632000 & 0.862000 & 1.000000 \\
    \bottomrule
    \end{tabular}}
    \caption{Ma trận tương quan giữa các biến quan trọng (theo hệ số OLS) và target.}
\end{table}

\subsubsection*{A3. Kết quả chọn lambda bằng k-fold CV}
\begin{table}[H]
    \centering
    \small
    \begin{tabular}{rrr}
    \toprule
    lambda & log10\_lambda & cv\_mse \\
    \midrule
    0.100000 & -1.000000 & 1.993500 \\
    0.046416 & -1.333333 & 1.993500 \\
    0.021544 & -1.666667 & 1.993500 \\
    0.010000 & -2.000000 & 1.993500 \\
    0.004642 & -2.333333 & 1.993501 \\
    \bottomrule
    \end{tabular}
    \caption{Top 5 lambda theo CV-MSE (k=5).}
\end{table}

\subsubsection*{A4. Top hệ số OLS sau chuẩn hóa}
\begin{table}[H]
    \centering
    \small
    \begin{tabular}{lr}
    \toprule
    Feature & Coefficient \\
    \midrule
    PT08.S2(NMHC) & 8.523142 \\
    PT08.S3(NOx) & 1.305243 \\
    T & -0.938744 \\
    AH & 0.859431 \\
    PT08.S1(CO) & 0.513167 \\
    PT08.S4(NO2) & -0.438952 \\
    RH & -0.428300 \\
    Hour & -0.240929 \\
    NOx(GT) & 0.176963 \\
    PT08.S5(O3) & -0.165408 \\
    \bottomrule
    \end{tabular}
    \caption{Top 10 hệ số OLS theo trị tuyệt đối (dữ liệu chuẩn hóa).}
\end{table}



\subsection*{Phụ lục B: GitHub Repository}
\addcontentsline{toc}{subsection}{Phụ lục B: GitHub Repository}

Mã nguồn của dự án được lưu trữ tại địa chỉ sau: 
\url{https://github.com/TNTruong196/Prj2_TUDTK_N1.git}

\centering
\textit{Hết.}
\end{document}